\chapter{Transacções conjuntas ({\tt TxTangle})}
\label{chapter:txtangle}

Uma transacção não conta histórias; deixa, contudo, relações. Mineiros,
reservas de mineração, mercados com garantia e bolsas produzem ritmos e
estruturas próprios, e esses hábitos podem alimentar heurísticas de grafos
mesmo quando os montantes e os verdadeiros membros dos anéis permanecem
ocultos (veja-se a secção 5.1 de \cite{AnalysisOfLinkability}). A álgebra cala
muito, mas não corrige por si só os costumes de quem a usa.

Este capítulo conserva a proposta \emph{TxTangle}, inspirada no CoinJoin de
Bitcoin \cite{coinjoin-wiki}: vários participantes constroem uma única
transacção e procuram esconder a partição das entradas e saídas entre si.
Pretende-se que um observador não consiga decidir quais entradas financiaram
quais saídas, nem quantos participantes reais estiveram presentes. Pretende-se
ainda que cada participante aprenda tão pouco quanto possível acerca dos
outros.

\paragraph{Estado da proposta.}
TxTangle não é uma funcionalidade do Monero, não é um tipo de transacção do
consenso e não está implementado nas carteiras de referência. O que se segue é
uma construção de investigação, originalmente escrita para a época de MLSAG e
Bulletproofs. Uma proposta anterior, MoJoin, atribuía a coordenação a um
administrador de confiança; essa confiança contrariava o objectivo de
fungibilidade e o trabalho não prosseguiu. No fim do capítulo distinguiremos a
ideia algébrica, que continua instrutiva, do protocolo histórico, que não pode
ser apresentado como software disponível.

\section{O núcleo algébrico}
\label{sec:building-txtangle}

Numa transacção RingCT, os pseudo-compromissos das entradas e os compromissos
das saídas satisfazem
\begin{align}
 \sum_j \widetilde C^a_j-
 \left(\sum_t C^b_t+fH\right)=0,
 \label{eq:txtangle-global-balance}
\end{align}
onde \(f\) é a taxa. Se cada participante construísse uma sub-transacção
separadamente equilibrada e todas fossem depois coladas, qualquer observador
poderia procurar subconjuntos que satisfizessem a mesma equação. O grande
segredo acabaria reduzido a um exercício de somas.

TxTangle propõe quebrar esses balanços parciais sem quebrar o balanço global.
Sejam \(N\) participantes. Para cada par não ordenado \(\{u,v\}\), escolhe-se
um escalar secreto \(s_{u,v}\in\mathbb Z_l\) e adopta-se a convenção
antissimétrica
\[
 s_{v,u}=-s_{u,v}.
\]
O participante \(u\) recebe o \emph{offset}
\[
 \delta_u=\sum_{v\ne u}s_{u,v}
\]
e adiciona \(\delta_uG\) à máscara de um dos seus pseudo-compromissos. Então
\begin{align}
 \sum_{u=1}^{N}\delta_u
 &=\sum_{\{u,v\}}(s_{u,v}+s_{v,u})=0.           \label{eq:txtangle-offset-zero}
\end{align}
Logo, a Equação~\eqref{eq:txtangle-global-balance} continua válida. Para um
subconjunto próprio \(S\subsetneq\{1,\ldots,N\}\), porém,
\begin{align}
 \sum_{u\in S}\delta_u
   =\sum_{\substack{u\in S\\v\notin S}}s_{u,v}, \label{eq:txtangle-cut}
\end{align}
porque os termos internos de \(S\) se cancelam aos pares. Sob a hipótese de
que pelo menos um termo de corte é desconhecido e uniforme, o lado direito é
uniforme em \(\mathbb Z_l\). Um balanço parcial deixa, portanto, um resíduo
aleatório; só a reunião de todos os participantes faz desaparecer a dívida.

É esta a parte verdadeiramente matemática de TxTangle. Os \emph{offsets}
alteram as máscaras, não os montantes, e cancelam-se apenas no agregado. Não
demonstram, todavia, que o protocolo distribuído que os produz seja seguro:
essa conclusão exige autenticação, vinculação do transcrito, resistência a
abortos e um modelo explícito do adversário.

\section{Construção descentralizada histórica}

\subsection{Canal de comunicação em grupo}
\label{subsec:n-way-channel}

O número de participantes reais não pode exceder o menor dos números de
entradas e saídas, pois cada participante precisa de contribuir com pelo menos
uma de cada. A proposta original atribui, além disso, uma identidade pseudónima
distinta a cada saída. Assim se forma um canal com \(n\) identidades sem se
publicar directamente quantas pessoas as controlam.

Uma sessão admite \(2\leq n\leq16\), limite que coincide com as dezasseis
saídas cobertas por uma Bulletproof+ nas regras actuais. A sala abre em
\(t_0\), fecha inscrições em \(t_1\) e fixa previamente a prioridade da taxa,
o método de estimação, a versão de transacção e os restantes parâmetros que
poderiam criar uma impressão digital. Em \(t_1\), cada identidade publica uma
chave pública. Dessas chaves derivam-se os segredos de canal e os escalares
\(s_{u,v}\).

As mensagens relativas a entradas eram autenticadas, na proposta, por
assinaturas SAG (secção \ref{SAG_section}); as relativas a saídas, por bLSAG
(secção \ref{blsag_note}) sobre o conjunto de identidades. Reutilizar a mesma
imagem de chave entre mensagens da mesma saída liga essas mensagens umas às
outras sem revelar qual identidade as produziu. É uma separação delicada:
ligar de menos permite substituições; ligar de mais recompõe o participante.

\subsection{Cinco rondas de mensagens}
\label{subsec:message-rounds-txtangle}

O protocolo histórico divide a construção em cinco rondas. Cada ronda tem um
intervalo próprio, e as publicações são espalhadas dentro dele para que a
proximidade temporal não denuncie grupos de entradas ou saídas.

\begin{enumerate}
 \item Cada identidade escolhe um escalar aleatório por saída, autentica-o com
 uma bLSAG e publica-o. A ordenação desses escalares determina os índices das
 saídas; o menor recebe \(t=0\). Publicam-se também, sob SAG, os tipos das
 entradas. Conhecidos os números e tipos de entradas e saídas, todos calculam
 o mesmo peso estimado e a mesma taxa.

 A divisão inteira da taxa é fixada por regra comum. O índice de saída mais
 baixo paga o resto. O campo \emph{extra}, a ordenação das saídas e todos os
 valores opcionais têm igualmente de seguir uma política comum: uma
 transacção conjunta que ostente um uniforme especial anuncia aquilo que
 desejava esconder.

 \item Para cada par de identidades deriva-se \(s_{u,v}\); a ordem canónica
 das chaves decide o sinal. Cada participante constrói os compromissos das
 suas saídas, os montantes cifrados e os pseudo-compromissos, aplicando
 \(\delta_uG\) a um destes últimos. A proposta original construía ainda a
 primeira parte de uma Bulletproof agregada. Todo o material de saída era
 autenticado por bLSAG; o material de entrada, por SAG. No fim da ronda, cada
 participante verifica o balanço global.

 \item Se o balanço for correcto, cada identidade calcula o primeiro desafio
 comum da Bulletproof e publica a contribuição seguinte para a prova. Uma
 divergência no transcrito deve abortar a sessão: continuar com desafios
 diferentes não produz uma prova parcialmente válida, mas duas provas
 inválidas inteiras.

 \item Fixam-se os membros dos anéis, as imagens de chave, os
 pseudo-compromissos, os endereços ocultos, as chaves públicas de transacção e
 a contribuição final para a prova de intervalo. A ordem e a autenticação das
 mensagens procuram impedir que o coordenador troque uma peça sem que o seu
 autor o detecte.

 \item Com o transcrito completo, todos obtêm a mesma prova de intervalo e a
 mesma mensagem a assinar. Cada participante produz as assinaturas das suas
 entradas e publica-as no canal. Quem reunir todas as componentes pode
 serializar e difundir a transacção.
\end{enumerate}

Estes passos descrevem MLSAG e Bulletproofs clássicos, não a construção activa
com CLSAG e Bulletproof+. A diferença não é uma troca de nomes: mudam os
desafios, os dados assinados, as condições de agregação e a gestão segura dos
\emph{nonces}.

\subsection{Chaves públicas de transacção e ataque de Janus}

Se todos conhecessem uma única chave privada de transacção \(r\), qualquer
participante poderia testar os endereços ocultos dos demais contra uma lista
de endereços conhecidos. A proposta usa, por isso, uma chave pública de
transacção distinta por saída, como acontece quando há destinatários em
sub-endereços (secção \ref{sec:subaddresses}).

O texto original considerava ainda uma chave pública `base' destinada à
mitigação de Janus \cite{janus-mitigation-issue-62}. Essa chave seria uma soma
de contribuições aleatórias, de modo que ninguém conhecesse sozinho o escalar
agregado. O cancelamento de chaves discutido na secção
\ref{subsec:drawbacks-naive-aggregation-cancellation} não autoriza aqui um
gasto, mas continua necessário vincular cada contribuição ao transcrito para
impedir substituições.

Ainda assim, quem recebe num sub-endereço pode reconhecer padrões: muitas
chaves públicas de transacção, muitas saídas ou uma chave base associável ao
troco. Há duas obrigações: cifrar os dados e disfarçar o formato. Cumprir a
primeira e esquecer a segunda seria uma economia bastante cara.

\subsection{Fraquezas}
\label{subsec:weaknesses-txtangle}

Há dois ataques elementares. Primeiro, um adversário pode controlar muitas
identidades e \emph{poluir} a sessão, reduzindo o conjunto honesto a um grupo
pequeno ou vazio \cite{coinjoin-pollution}. Sem reputação, credenciais escassas
ou outro custo de admissão, uma identidade não prova que atrás dela exista uma
pessoa distinta.

Segundo, o adversário pode abortar repetidamente e comparar as tentativas
seguintes. As chaves públicas de transacção, máscaras, escalares das provas e
\emph{nonces} de assinatura têm de ser regenerados; reutilizá-los pode revelar
segredos ou ligar transcritos. Os membros fictícios dos anéis, pelo contrário,
devem permanecer fixos quando se reutiliza a mesma entrada, pois a intersecção
de anéis entre tentativas pode denunciar o membro verdadeiro. A regra mais
segura é não reutilizar a entrada após abortos, embora isso imponha custos de
liquidez e selecção de moedas.

Acrescem a negação de serviço, a estimação inconsistente da taxa, mensagens
equívocas, contribuições inválidas para a prova agregada e a fuga de metadados
na rede. Nenhuma equação de cancelamento resolve esses problemas.

\section{TxTangle servido}
\label{sec:hosted-txtangle}

No modelo descentralizado resta saber quem abre as salas, faz avançar as
rondas e distribui as mensagens. A solução mais directa introduz um anfitrião.
Ele simplifica a disponibilidade e o relógio, mas ocupa uma cadeira demasiado
boa: vê ligações, tempos e volumes, e pode apresentar listas de participantes
diferentes a pessoas diferentes.

\subsection{Comunicação através de I2P}
\label{subsec:txtangle-host-communication}

A proposta sugere I2P para dissociar as mensagens das ligações de rede
\cite{exa-blockchain-analysis}. Túneis distintos transportariam pedidos de
inscrição, mensagens de entrada e mensagens de saída. Essa disciplina reduz
correlações triviais, mas não transforma o anfitrião em observador cego: a
temporização, o tamanho, a reutilização de circuitos e a possibilidade de ele
controlar pontos de entrada continuam no modelo de ameaça.

O serviço executaria quatro tarefas:

\begin{enumerate}
 \item \emph{Formar sessões.} Cada candidato declara o número de saídas, ou
 entrega as chaves das identidades pseudónimas, e o anfitrião reúne uma sala
 compatível. Aprende assim o volume declarado por ligação, embora não deva
 aprender a associação final entre entradas e saídas.

 \item \emph{Transportar rondas.} O anfitrião recolhe as mensagens durante o
 intervalo, publica o conjunto canónico no fim e aguarda uma margem antes da
 ronda seguinte. Todos devem comprometer-se com o mesmo \emph{hash} do
 transcrito.

 \item \emph{Separar circuitos.} Mensagens ligadas à mesma saída podem usar o
 mesmo túnel; mensagens de saídas ou entradas diferentes não o devem fazer.
 Pedidos fictícios podem reduzir a correlação entre consultas e inscrições,
 mas aumentam custo e não constituem uma prova de anonimato.

 \item \emph{Impedir a interposição.} Se o anfitrião entregar a Alice uma
 lista em que todas as outras identidades são suas, aprende a partição de
 Alice. A defesa proposta vincula as chaves públicas de transacção ao conjunto
 canónico \(\mathbb S_{\rm id}\):
 \[
  R_t=
  \mathcal H_n(T_{\rm agg},\mathbb S_{\rm id},R_t^{(0)})\,R_t^{(0)},
  \qquad R_t^{(0)}=r_tG.
 \]
 Participantes que receberam listas diferentes obtêm chaves diferentes e não
 conseguem concluir a mesma transacção. Numa implementação moderna, a
 vinculação teria de abranger a versão, a sessão, todas as mensagens e a ordem
 canónica, e não apenas esta chave.
\end{enumerate}

\subsection{Operação como serviço}
\label{subsec:txtangle-host-service}

Um anfitrião pode cobrar sem criar contas: acrescenta uma saída sua e exige
que os participantes a financiem. A taxa de serviço e a taxa de mineiro são
divididas por uma regra determinística; uma identidade canonicamente escolhida
paga os restos das divisões inteiras.

Como o anfitrião não fornece uma entrada, não possui um pseudo-compromisso que
cancele a máscara da sua saída. A proposta reparte essa máscara por escalares
destinados aos participantes. Se \(h\) é a máscara da saída do anfitrião, as
contribuições \(\eta_u\) têm de satisfazer
\[
 \sum_u\eta_u=-h,
\]
e cada participante incorpora a sua parcela num pseudo-compromisso. O sinal
depende da convenção adoptada na Equação~\eqref{eq:txtangle-global-balance}; o
invariante indispensável é que a soma final das máscaras seja zero.

O anfitrião pode abortar perante mensagens inválidas ou falta de pagamento e,
na ronda final, difundir a transacção e o seu identificador. Esta faculdade é
útil e perigosa: oferece disponibilidade quando tudo corre bem e um oráculo de
abortos quando não corre.

\section{Construção com administrador de confiança}
\label{sec:dealer-txtangle}

O modelo inspirado em MoJoin abandona a tentativa de esconder a partição do
administrador. Este recebe os números e tipos de entradas e saídas, escolhe os
índices, calcula a taxa e distribui os \emph{offsets}. A execução histórica é:

\begin{enumerate}
 \item Para cada par de participantes, o administrador escolhe um escalar com
 sinais opostos. Envia privadamente a cada participante a soma que lhe cabe,
 a fracção da taxa e os índices das suas saídas.

 \item Cada participante constrói a sua sub-transacção: endereços ocultos,
 compromissos, montantes cifrados, pseudo-compromissos e dados dos anéis.
 Adiciona o \emph{offset} recebido a um pseudo-compromisso e envia ao
 administrador a primeira contribuição para a prova de intervalo.

 \item O administrador distribui o conjunto canónico dessas contribuições;
 cada participante calcula o desafio seguinte e devolve a segunda parte.

 \item Depois da última contribuição, o administrador conclui a prova de
 intervalo e envia a mensagem final a assinar.

 \item Cada participante assina as suas entradas. O administrador reúne as
 assinaturas, serializa a transacção e difunde-a.
\end{enumerate}

O protocolo fica mais simples porque o administrador conhece aquilo que o
modelo descentralizado tentava esconder. A privacidade passa então de uma
propriedade criptográfica para uma promessa operacional.

\section{Avaliação face ao Monero actual}
\label{sec:txtangle-modern-status}

O código-fonte de referência foi verificado na revisão de 14 de Agosto de
2026:
\begin{center}
 \small\texttt{3646f648db57f60cca86430e25a635d19fa9b92a}.
\end{center}
Não contém TxTangle, MoJoin nem uma interface CoinJoin para Monero.
\footnote{A construção ordinária está em
\texttt{src/wallet/wallet2.cpp} e \texttt{src/ringct/rctSigs.cpp}; a construção
distribuída existente pertence às contas de multi-assinatura e usa
\texttt{src/multisig/multisig\_tx\_builder\_ringct.cpp}. Nenhum destes caminhos
implementa o protocolo deste capítulo.}

Há, pois, três afirmações que não devemos confundir:

\begin{enumerate}
 \item \emph{A álgebra é plausível.} Os \emph{offsets} antissimétricos das
 Equações~\eqref{eq:txtangle-offset-zero}--\eqref{eq:txtangle-cut} preservam o
 balanço global e escondem balanços próprios de subconjuntos.

 \item \emph{Uma transacção final poderia usar o formato ordinário.} Se todos
 os compromissos, provas e assinaturas fossem válidos, o consenso não
 precisaria de uma etiqueta `TxTangle': verificaria apenas uma transacção
 RingCT. Isto torna concebível uma implementação na camada das carteiras e da
 coordenação; não significa que ela exista.

 \item \emph{O transcrito histórico não é uma especificação actual.} O
 caminho activo usa CLSAG e Bulletproof+, com etiquetas de vista e
 regras modernas para o campo \emph{extra}. FCMP++ e Carrot alteram novamente
 a autorização, a pertença e a construção das saídas, mas não estavam activos
 no consenso na revisão considerada (capítulo \ref{chapter:fcmp-plus-plus}).
\end{enumerate}

Uma implementação séria teria de fornecer, no mínimo, uma especificação
canónica do transcrito; provas de posse das entradas; geração distribuída de
Bulletproof+ ou outra forma segura de agregar provas; protecção dos
\emph{nonces}; validação de cada mensagem antes da ronda seguinte; regras de
aborto e recuperação; defesa contra Sybil e poluição; testes diferenciais com
o verificador do consenso; análise de metadados; vectores de teste; e uma
auditoria independente. Teria ainda de provar que um participante malicioso
não consegue roubar fundos, criar inflação, ligar honestos ou obter assinaturas
sobre uma transacção diferente.

TxTangle é, em suma, uma proposta matemática real e uma implementação ausente.
O cancelamento dos \emph{offsets} merece ficar; a alegação de que o protocolo
vive hoje em XMR, essa deve ficar de fora. Entre uma equação elegante e uma
carteira segura há sempre um corredor comprido --- e, neste caso, ainda não foi
construída a porta.
